\SourceLeaf{019}per omnes valores reales ipsarum \(x,x',x''\) etc. extensa, puta ab \(x=-\infty\) usque ad \(x=+\infty\), ab \(x'=-\infty\) usque ad \(x'=+\infty\) etc.

\Article{14}
Reducta itaque functione \(y\) ad formam aggregati talium partium
\[
Ax^\alpha x'{}^\beta x''{}^\gamma\mathord{.}\,\mathord{.}\,\mathord{.}\,\mathord{.}
\]
valor integralis \(\int y\psi y.\mathrm{d}y\) per omnes valores ipsius \(y\) extensi, seu valor medius ipsius \(y\), aequalis erit aggregato partium
\[
A\times\int x^\alpha\varphi x.\mathrm{d}x\times\int x'{}^\beta\varphi x'.\mathrm{d}x'\times\int x''{}^\gamma\varphi x''.\mathrm{d}x''\mathord{.}\,\mathord{.}\,\mathord{.}\,\mathord{.}
\]
ubi integrationes extendendae sunt ab \(x=-\infty\) usque ad \(x=+\infty\), ab \(x'=-\infty\) usque ad \(x'=+\infty\) etc.; sive quod eodem redit, aggregato partium quae oriuntur, dum pro singulis potestatibus \(x^\alpha,x'{}^\beta,x''{}^\gamma\) etc. ipsarum valores medii substituuntur, cuius theorematis gravissimi veritas etiam ex aliis considerationibus facile derivari potuisset.

\Article{15}
Applicemus ea, quae in art. praec. exposuimus, ad casum specialem, ubi
\[
y=\frac{xx+x'x'+x''x''+\text{etc.}}{\sigma}
\]
denotante \(\sigma\) multitudinem partium in numeratore. Valor medius ipsius \(y\) hic illico invenitur \(=mm\), accipiendo characterem \(m\) in eadem significatione ac supra. Valor verus quidem ipsius \(y\) in casu determinato maior minorve evadere potest medio, perinde ac valor verus termini simplicis \(xx\): sed probabilitas quod valor fortuitus ipsius \(y\) a medio \(mm\) haud sensibiliter aberret, continuo magis ad certitudinem appropinquabit crescente multitudine \(\sigma\). Quod quo clarius eluceat, quum probabilitatem ipsam exacte determinare non sit in potestate, investigemus errorem medium metuendum, dum supponimus \(y=mm\). Manifesto per principia art. 6 hic error erit radix quadrata valoris medii functionis
\[
\left(\frac{xx+x'x'+x''x''+\text{etc.}}{\sigma}-mm\right)^2
\]
ad quem eruendum sufficit observare, valorem medium termini talis \(\frac{x^4}{\sigma\sigma}\) esse \(=\frac{n^4}{\sigma\sigma}\) (utendo charactere \(n\) in significatione art. 11), valorem medium autem termini \SourceLeaf{020}talis \(\frac{2xxx'x'}{\sigma\sigma}\) fieri \(=\frac{2m^4}{\sigma\sigma}\), unde facillime deducitur valor medius istius functionis
\[
=\frac{n^4-m^4}{\sigma}
\]

Hinc discimus, si copia satis magna errorum fortuitorum ab invicem independentium \(x,x',x''\) etc. in promtu sit, magna certitudine inde peti posse valorem approximatum ipsius \(m\) per formulam
\[
m=\BareRoot{\frac{(xx+x'x'+x''x''+\text{etc.})}{\sigma}}
\]
erroremque medium in hac determinatione metuendum, respectu quadrati \(mm\), esse
\[
=\BareRoot{\frac{n^4-m^4}{\sigma}}
\]
Ceterum, quum posterior formula implicet quantitatem \(n\), si id tantum agitur, ut idea qualiscunque de gradu praecisionis istius determinationis formari possit, sufficiet, aliquam hypothesin respectu functionis \(\varphi\) amplecti. E. g. in hypothesi tertia art. 9, 11 iste error fit \(=mm\BareRoot{\frac2\sigma}\). Quod si minus arridet, valor approximatus ipsius \(n^4\) ex ipsis erroribus adiumento formulae
\[
\frac{x^4+x'{}^4+x''{}^4+\text{etc.}}{\sigma}
\]
peti poterit. Generaliter autem affirmare possumus, praecisionem duplicatam in ista determinatione requirere errorum copiam quadruplicatam, sive pondus determinationis ipsi multitudini \(\sigma\) esse proportionale.

Prorsus simili modo, si observationum errores partem constantem involvunt, huius valor approximatus eo tutius e medio arithmetico multorum errorum colligi poterit, quo maior horum multitudo fuerit. Et quidem error medius in hac determinatione metuendus exprimetur per
\[
\BareRoot{\frac{mm-kk}{\sigma}}
\]
si \(k\) designat partem constantem ipsam atque \(m\) errorem medium observationum parte constante nondum purgatarum, sive simpliciter per \(\frac{m}{\BareRoot{\sigma}}\), si \(m\) denotat errorem medium observationum a parte constante liberatarum (v. art. 8).

\Article{16}
In artt. 12---15 supposuimus, errores \(x,x',x''\) etc. ad idem observationum genus pertinere, ita ut singulorum probabilitates per eandem functionem expri\SourceLeaf{021}mantur. Sed sponte patet, disquisitionem generalem artt. 12---14 aeque facile ad casum generaliorem extendi, ubi probabilitates errorum \(x,x',x''\) etc. per functiones diversas \(\varphi x,\varphi'x',\varphi''x''\) etc. exprimantur, i. e. ubi errores illi pertineant ad observationes praecisionis seu incertitudinis diversae. Supponamus, \(x\) esse errorem observationis talis, cuius error medius metuendus sit \(=m\); nec non \(x',x''\) etc. esse errores aliarum observationum, quarum errores medii metuendi resp. sint \(m',m''\) etc. Tunc valor medius aggregati \(xx+x'x'+x''x''+\text{etc.}\) erit \(mm+m'm'+m''m''+\text{etc.}\) Iam si aliunde constat, quantitates \(m,m',m''\) etc. esse in ratione data, puta numeris \(1,\mu',\mu''\) etc. resp. proportionales, valor medius expressionis
\[
\frac{xx+x'x'+x''x''+\text{etc.}}{1+\mu'\mu'+\mu''\mu''+\text{etc.}}
\]
erit \(=mm\). Si vero valorem eiusdem expressionis determinatum, prout fors errores \(x,x',x''\) etc. offert, ipsi \(mm\) aequalem ponimus, error medius, cui haec determinatio obnoxia manet, simili ratione ut in art. praec. invenitur
\[
=\frac{\BareRoot{(n^4+n'{}^4+n''{}^4+\text{etc.}-m^4-m'{}^4-m''{}^4-\text{etc.})}}{1+\mu'\mu'+\mu''\mu''+\text{etc.}}
\]
ubi \(n',n''\) etc. respectu observationum, ad quas pertinent errores \(x',x''\) etc., idem denotare supponuntur, atque \(n\) respectu observationis primae. Quodsi itaque numeros \(n,n',n''\) etc. ipsis \(m,m',m''\) etc. proportionales supponere licet, error ille metuendus medius fit
\[
=\frac{\BareRoot{(n^4-m^4)}.\BareRoot{(1+\mu'{}^4+\mu''{}^4+\text{etc.})}}{1+\mu'\mu'+\mu''\mu''+\text{etc.}}
\]

At haecce ratio, valorem approximatum ipsius \(m\) determinandi non est ea, quae maxime ad rem facit. Quod quo clarius ostendamus, consideremus expressionem generaliorem
\[
y=\frac{xx+\alpha'x'x'+\alpha''x''x''+\text{etc.}}{1+\alpha'\mu'\mu'+\alpha''\mu''\mu''+\text{etc.}}
\]
cuius valor medius quoque erit \(=mm\), quomodocunque eligantur coëfficientes \(\alpha',\alpha''\) etc. Error autem medius metuendus, dum valorem determinatum ipsius \(y\), prout fors errores \(x,x',x''\) etc. offert, ipsi \(mm\) aequalem supponimus, invenitur per principia supra tradita
\[
=\frac{\BareRoot{(n^4-m^4+\alpha'\alpha'(n'{}^4-m'{}^4)+\alpha''\alpha''(n''{}^4-m''{}^4)+\text{etc.})}}{1+\alpha'\mu'\mu'+\alpha''\mu''\mu''+\text{etc.}}
\]
